\documentclass[11pt,a4paper]{article}
\usepackage[margin=25mm]{geometry}
\usepackage[T1]{fontenc}
\usepackage{lmodern,microtype,amsmath,amssymb,amsthm,booktabs,graphicx}
\usepackage{xcolor,enumitem,caption}
\usepackage[colorlinks=true,linkcolor=blue!40!black,citecolor=blue!40!black,urlcolor=blue!40!black]{hyperref}
\hypersetup{pdftitle={Planning Mine Escape Routes with Zone-Based Flood Alarms},pdfauthor={Jingchao Wang}}
\captionsetup{font=small,labelfont=bf}
\setlist{itemsep=2pt,topsep=4pt,leftmargin=1.6em}
\setlength{\emergencystretch}{2em}
\newtheorem{proposition}{Proposition}
\newcommand{\Sc}{\mathcal{S}}
\title{Planning Mine Escape Routes\\with Zone-Based Flood Alarms}
\author{Jingchao Wang}
\date{}
\begin{document}
\maketitle
\begin{abstract}
Mine escape routes are fixed in advance, usually as the shortest way to a shaft, but when water breaks in, the alarm reports only which zone of the mine is flooding. We ask what a fixed escape plan needs to know about where the water enters. On a 248-km tunnel network from a real coal mine, we simulate a flood from each of its 1,587 junctions and compare three plans: the shortest route, a simple rule that avoids the alarm zone, and the shortest route that is guaranteed for every possible source in the zone, which we compute exactly. Three findings stand out. First, shortest routes fail not because the warning is late but because they converge on the shafts, where the water arrives first; with four alarm zones and a 30-minute delay, they can be guaranteed for only 59\% of (worker, flood) pairs. Second, the shortest route that survives every possible source in the zone raises this to 79\%; a simple rule that steers clear of the alarm zone, without checking the flood, reaches 69\% but with much longer walks. Third, almost all of the remaining gap to the 99.8\% possible with the exact source is a lack of information, not of time: finer alarm zones close it, and drawing zone boundaries through the shafts is worth about as much as doubling the number of zones.
\end{abstract}

\section{Introduction}
A water inrush happens when groundwater, or water stored in old workings, suddenly breaks into an operating mine. It is among the most dangerous accidents underground \cite{cumcm,sun2016,sun2024}, because the water and the miners move through the same tunnels. A large mine can have hundreds of kilometers of tunnels connected to the surface by only a few shafts. Water entering at one point runs downhill, fills low sections, and can cut off passages far from where it started. Workers who are still underground have to reach a shaft before the tunnels on their way become too deep to pass.

Escape routes are therefore planned in advance. Each work area is assigned a route to a shaft, the route is marked underground, and workers practice it in drills \cite{nrc2013}. In an emergency people follow what they have practiced; nobody computes a fresh route for each worker. The natural route to practice is the shortest one.

The difficulty is that a route fixed in advance cannot know where the water will come from. Even after the alarm sounds, the source is usually known only roughly. A water-level sensor or a report from a work face tells the control room that water has appeared in some part of the mine, and the alarm passes that zone on to workers. A route that is safe for a flood in one corner of that zone may run straight into a flood from another corner. The warning may also arrive late, because water can travel some distance before it reaches a sensor.

Earlier work has treated these two sides separately. Escape-route studies find the fastest path through a mine whose flooding has been simulated in advance \cite{wu2020}, and sensor studies place water-level sensors so that an inrush is detected quickly \cite{wu2021}. What connects them is the plan itself: what route should a worker be told to take when the alarm reports only a zone? This paper asks what a fixed escape plan needs to know about where the water comes from, and what it should do with the little that an alarm reveals. We break the question into three parts.
\begin{enumerate}
\item How often does the shortest route fail when only the zone is known, and why?
\item How much better can a route do if it is checked against every possible source in the zone, how much walking does it cost, and does a simple rule do as well?
\item When no fixed route works, is what is missing better information about the source, or an earlier warning?
\end{enumerate}

The questions grew out of a recent modeling competition. Problem~D of the 2025 China Undergraduate Mathematical Contest in Modeling asked teams to model how inrush water spreads through a mine's tunnels and to plan escape routes for the miners \cite{cumcm}. In that problem the inrush point is given exactly and the warning goes out one minute after the inrush begins. A real alarm rarely offers either, and this paper asks what changes when it does not. Our flood model keeps the competition's rules almost unchanged; because those rules were set for a contest rather than measured in a mine, the model is uncalibrated, and we read its results as a test case for the planning method.

We study these questions on a tunnel network taken from the plan of a real coal mine, with 1,592 junctions and about 248\,km of tunnels. For each of its 1,587 junctions we simulate a flood that starts there, using a simple model of how water spreads and when each tunnel becomes impassable. We then compare plans for every starting point and every alarm zone. Besides the shortest route, we test a rule that a planner could apply with a map alone: avoid the tunnels of the alarm zone. Against both we set the \emph{robust route}, the shortest route that reaches a shaft for every possible source in the zone. Finding it means searching for one route that works in hundreds of floods at once, which we do exactly. Finally, we vary the alarm itself, from no location information at all to 128 small zones, and compare the result with the ceiling set by knowing the exact source.

The main findings are as follows. With four alarm zones and a 30-minute warning delay, shortest routes can be guaranteed for 59\% of all (worker, flood) pairs, the avoid-the-zone rule for 69\%, and robust routes for 79\%, whereas knowing the exact source would allow 99.8\%. The shortest routes do not fail because the warning is late: most fail even for a worker who leaves at once, and they are cut within a few hundred meters of the shafts, where routes from all over the mine meet. The avoid-the-zone rule fixes many of the same workers, but often by sending them on very long detours that a robust route avoids. Nearly all of the remaining gap is a matter of information rather than time, and it closes as the alarm zones become finer. Where the zone boundaries fall matters as much as how many zones there are. A zone that contains the whole shaft area is a weak point that more zones elsewhere do not remove, and two zones divided at the shafts protect 52\% of pairs, against 31--36\% for two zones of equal size.

The case study quantifies how source-location information changes the coverage and walking cost of fixed escape plans.

\paragraph{Related work.}
Wu et al.\ \cite{wu2020} find earliest-arrival escape paths through a mine network whose roadways flood over time, using the results of a flood simulation for a given inrush, and Zhang et al.\ \cite{zhang2022} model the spreading itself in more detail. Water-level sensors can be placed to detect an inrush quickly \cite{wu2021}; our results suggest that where zone boundaries fall relative to the shafts also matters for what the detection can be used for. Evacuation planning on networks more generally has a large optimization literature \cite{bayram2016}.

Our computational tools come from time-dependent and robust routing. Shortest paths with travel times that change over time go back to Dreyfus \cite{dreyfus1969}; when arriving earlier never means leaving later, arrival times can be compared directly and label-setting search stays exact \cite{kaufman,orda}. Comparing partial routes on several criteria at once follows multicriteria shortest-path methods \cite{martins1984}, including time-dependent versions \cite{hamacher2006}. Uncertain travel times have been handled through expected arrival times \cite{millerhooks2000} and through robust shortest paths that guard against the worst of several scenarios \cite{yu1998,bertsimas2003}; Ben-Tal et al.\ \cite{bental2009} survey robust optimization in general. Our search adds a scenario only when a candidate route fails it, in the spirit of cutting-set methods \cite{mutapcic2009}. Choosing among prepared plans after part of the uncertainty is revealed is studied as $K$-adaptability \cite{hanasusanto}; here the revealed part is the alarm zone, and we study how its resolution changes what a plan can guarantee.

Section~\ref{sec:setup} sets up the problem and Section~\ref{sec:compute} explains how robust routes are computed. Sections~\ref{sec:why}--\ref{sec:info} answer the three questions in turn, Section~\ref{sec:robust} tests how the answers change with the inflow and the set of possible sources, and Section~\ref{sec:limits} discusses limitations.

\section{The planning problem}\label{sec:setup}
\subsection{The mine}
Our test case is a tunnel network extracted from an anonymized CAD plan of a working longwall coal mine (Figure~\ref{fig:geometry}). It has 1,592 junctions and 2,084 tunnels with a total length of about 248\,km. A main roadway runs diagonally across the mine, and longwall panels branch off it on both sides. The seam dips gently: tunnel floors in the west lie about 80\,m higher than those in the far northeast, so water tends to run east and northeast. All four shafts stand on the main roadway, within about half a kilometer of each other near the middle of the mine. Once a worker reaches the bottom of a shaft we count the escape as complete; the trip up the shaft is not modeled.

To build the network we joined nearly touching lines on the plan, split lines where they cross, and connected the shafts to the tunnels. The plan gives roof elevations, which we use as relative floor elevations by assuming all tunnels have the same height.

\begin{figure}[tbp]
\centering\includegraphics[width=\linewidth]{figures/geometry3d.pdf}
\caption{The mine in three dimensions, with the vertical scale exaggerated 60 times. Dark lines are the tunnels at their floor elevation; light gray lines are the same tunnels projected onto the floor of the box. The workings dip gently from the longwall panels in the west, about 80\,m above the lowest point, toward the far northeast. The four shafts (red) stand close together on the main roadway near the middle of the mine; their height above the workings is not to scale.}\label{fig:geometry}
\end{figure}

\subsection{Sources and alarm zones}
Any junction except the four shafts and one auxiliary node can be the place where water enters, which we call the \emph{source}. That gives 1,587 possible sources, and the same 1,587 junctions are the possible starting points of workers. Each flood has exactly one source, and we treat all sources as equally likely.

The alarm divides the mine into zones and reports the zone that contains the source. It does not report the source itself. We use two kinds of zones. The first is the four \emph{quadrants} formed by the median $x$ and $y$ coordinates of the junctions (Figure~\ref{fig:plan}); they are the simplest possible split. Their boundaries happen to cross right at the shaft cluster, a fact that turns out to matter. The second is a family of finer zones that we use to measure the value of better information. We cut the mine along a straight line into two halves that contain the same tunnel length, cut each half along the perpendicular direction, and so on, which gives 2, 4, 8, \ldots, 128 zones of roughly equal tunnel length. With 128 zones, a zone holds about 2\,km of tunnel. To check that the results do not depend on where the cuts fall, we repeat the construction with the first cut at three angles, 0, 60 and 120 degrees. In a real mine, the zones would follow the placement of water sensors; ours are geometric stand-ins that let us vary the resolution in a controlled way.

\begin{figure}[tbp]
\centering\includegraphics[width=.62\linewidth]{figures/plan.pdf}
\caption{Plan view of the mine and the four quadrant alarm zones (dashed boundaries). All four shafts (red squares) sit on the main roadway where the zone boundaries cross; the inset enlarges this area.}\label{fig:plan}
\end{figure}

\subsection{How late the warning can be}
The alarm also comes late. Water may travel some distance before it reaches a sensor, and passing the warning on to workers takes more time. If the inrush starts at time $t_0$ and the worker leaves on receiving the alarm at time $t_{\rm rec}$, the worker sets off a delay $d=t_{\rm rec}-t_0$ after the inrush began. The planner cannot know $d$ in advance, so we ask for routes that work for every delay up to a target $\tau$. Our main target is 30 minutes, and we vary it from zero to two hours. The target is about how late the worker may set off, not about how long the walk takes.

\subsection{The flood and walking model}\label{sec:hazard}
To check a route, we need to know, for each source, when each tunnel first gets wet and when it becomes too deep to pass. We compute these times with an event-driven model that tracks the volume of water in every tunnel, which follows the rules of the competition problem described in the introduction \cite{cumcm}. More detailed models add pipe-flow dynamics, energy losses at junctions and pumping \cite{zhang2022}; we use the simpler model because the routing questions need flood times for all 1,587 sources. Water arriving at a junction divides equally among the level or downhill tunnels leaving it, and goes uphill only when no such tunnel is available. A thin layer of water covers a tunnel before the water moves on beyond it. Dead-end branches fill up, and water then backs up behind them. The depth in a tunnel is its stored volume divided by its floor area. Appendix~\ref{app:flood} gives the full rules.

The main runs use the competition's inflow of 30\,m$^3$/min ($0.5\,\mathrm{m^3/s}$), its tunnel height of 3\,m, and its 0.1\,m water layer at the advancing front. A tunnel closes for good when its mean depth reaches 0.3\,m, the depth above which the competition advises against wading. We depart from the competition in two places. The tunnels are 4.99\,m wide, as estimated from the mine plan, rather than 4\,m. And workers walk more slowly than the competition's 4, 2 and 1\,m/s: 1.3\,m/s in dry tunnels, and in wet tunnels 0.7\,m/s with the flow and 0.5\,m/s against it. If water reaches a tunnel while a worker is inside it, the worker slows down from that moment. These values define a test case; they are not measured walking speeds or a validated safety threshold.

One property of this model simplifies everything that follows. Water never recedes, so a worker who enters a tunnel earlier never leaves it later. It follows that a route that works for a worker who sets off at $\tau$ also works for any earlier departure, and checking a whole range of delays only requires checking the latest one (Appendix~\ref{app:math}).

\subsection{Four plans}
A plan assigns each starting point $v$ and each alarm zone $Z$ one route, fixed in advance. Every route in every plan ends at the bottom of one of the four shafts; the plans differ in which shaft they choose and how they get there. We compare four plans.
\begin{description}[leftmargin=1.2em,itemsep=3pt]
\item[Shortest route.] Every worker takes the shortest route to the nearest shaft, whatever the alarm says. This is the natural default.
\item[Avoid-the-zone rule.] Among all routes to a shaft, the worker takes the shortest of those that spend the least distance in tunnels touching the alarm zone. A planner could apply this rule with a map alone. The rule does not look at the flood when it picks a route; we check its routes afterwards, with the same test as every other plan.
\item[Robust route.] Among all routes to a shaft, the shortest one that gets the worker out for every possible source in the zone, checked one flood at a time, and for every delay up to $\tau$. Unlike the avoid-the-zone rule, it uses the flood model to pick the route, and it may cross tunnels of the zone that the water reaches late. Finding it requires the search described in Section~\ref{sec:compute}. When no such route exists, the starting point cannot be protected by any fixed plan for that zone.
\item[Exact source.] For comparison only, the best route for a worker who is told the exact source. No alarm system delivers this, but it marks a ceiling that no plan can beat.
\end{description}
Let $\Sc_Z$ be the sources in zone $Z$ and $L(P)$ the length of route $P$. The robust route has length
\begin{equation}
\begin{split}
 L_Z^*(v,\tau)=\min\bigl\{L(P):{}&P\text{ leads from }v\text{ to a shaft and works}\\
 &\text{for departure at }\tau\text{ for every source in }\Sc_Z\bigr\},
\end{split}
 \label{eq:main}
\end{equation}
which is infinite when no route qualifies. Because it is the shortest route with the guarantee, its extra length over the shortest dry route measures the least walking that the guarantee demands, and any cap on walking can be applied afterwards.

\subsection{How we measure protection}
We say a plan \emph{protects} a starting point for a zone when the route it assigns works for every source in the zone and every delay up to $\tau$. For the quadrants, we count protected starting points outside and inside each zone separately. Workers inside the zone may be standing next to the source, and their prospects are much worse.

To compare zones of different sizes, we also use a whole-mine measure. Consider every pair of a starting point and a source, 1,587 by 1,587 pairs in all. A pair is protected if the plan's route for the source's zone is guaranteed for that starting point. The share of protected pairs is the chance that a randomly placed worker is protected in a flood from a randomly placed source. Each unprotected pair falls into one of two groups. If even the exact source would not help, because the worker has no way out of that flood after a delay of $\tau$, we call the pair \emph{no way out}. Otherwise, a route exists for that flood, but no single route works for the whole zone; we call the pair \emph{information-limited}, because a finer alarm could help it.

\section{Finding a route that works for every source}\label{sec:compute}
Computing the best route for each source separately and comparing the answers does not give one route that works for all of them, so the search has to consider all floods in the zone at once. We do this exactly for a finite set of sources.

\paragraph{One clock for each flood.}
Imagine that a partial route carries one clock for each flood, showing when the worker would reach the current junction in that flood. The search stores \emph{labels} at each junction, each recording a route length and the arrival times in all floods. Extending a route along a tunnel updates every clock, using that flood's wetting time, closure time and walking speed. If the tunnel would already be closed in any flood when the worker gets there, the extension is discarded. If one label at a junction is no longer than another and arrives no later in every flood, the other label can be dropped, since arriving earlier never hurts; this is the dominance rule of multicriteria shortest-path search \cite{martins1984,hamacher2006}. The search expands labels in order of length plus the remaining dry distance to a shaft, so the first complete route it finds is the shortest route that works in every flood.

\paragraph{Adding floods only when needed.}
A zone can hold hundreds of sources, and most of them do not matter for a given starting point. We start with the source that is worst for the shortest route, find the shortest route that works for it, and check that route against every source in the zone with a fast backward calculation. If some source defeats the route, we add that source and search again; otherwise the route is optimal for the whole zone. If no route works even for the sources added so far, none works for the zone. This is a cutting-set scheme \cite{mutapcic2009} with the flood model as the oracle that finds a defeating source. The procedure is exact and ends after at most as many rounds as there are sources (Appendix~\ref{app:math}).

\paragraph{Scale.}
Our experiments solve this problem 2,558,244 times, once for every starting point, zone, delay target and inflow considered. With time limits of three seconds per search and 30 seconds per query, all but seven finished. We re-solved those seven, all for workers inside zones that cover at least a quarter of the mine, with a limit of one hour. Each finished within 18 seconds and showed that no common route exists, so every result reported here is exact. We checked the implementation against brute-force enumeration on small networks, confirmed that the results obey the logical relations that must hold between the four plans, and walked a sample of routes forward against every source with separate code (Appendix~\ref{app:checks}).


\section{Why the shortest route fails}\label{sec:why}
We begin with the four quadrant zones and a 30-minute delay target. Table~\ref{tab:plans} gives, for each zone, how many starting points outside it each plan protects. Workers inside the zone are discussed at the end of Section~\ref{sec:fix}.

The shortest route works poorly for alarms in the two western zones. For a northwest alarm it protects 720 of the 1,243 starting points outside the zone (58\%), and for a southwest alarm 659 of 1,139 (58\%). The southeast and northeast zones lose far fewer: 45 and 63 starting points.

These failures are not caused by the late warning. Of the 523 northwest failures, 473 would happen even if the worker left at the moment the inrush began, and the same holds for 441 of the 480 southwest failures (Figure~\ref{fig:failures}b). Across all four zones, 91\% of the shortest routes that fail would fail without any delay at all.

They fail because of where the shafts are. To see where, we followed each failing shortest route through the flood that defeats it and recorded where the worker is stopped. For a northwest alarm, the median stopping point lies 130\,m along the tunnels from a shaft, and 508 of the 523 lie within a kilometer of one; for a southwest alarm the median is 280\,m. All four shafts stand close together on the main roadway, right where the northwest and southwest zones meet (Figure~\ref{fig:plan}). Water from a source in either zone runs into the junctions around the shafts, where routes from every part of the mine converge. A worker far from the flood, walking the shortest route, walks into it at the very end.

\begin{figure}[tbp]
\centering\includegraphics[width=\linewidth]{figures/failures.pdf}
\caption{Where and why shortest routes fail, at a 30-minute delay target. (a) Northwest alarm. Purple dots are the possible sources; colored dots are starting points outside the zone. Starting points whose shortest route fails (orange and red) cover much of the mine, because their routes are cut near the shafts (red squares), which lie at the edge of the zone. (b) For every quadrant zone, the share of outside starting points whose shortest route works, fails even if the worker leaves at once, or fails only because of the delay.}\label{fig:failures}
\end{figure}

Seen over the whole mine, a shortest-route plan with quadrant alarms can guarantee an escape for 58.7\% of all (worker, flood) pairs. The guarantee is a demanding standard, because it must hold for every source in the zone. A complementary measure asks how often the route actually fails for the source that occurs. By this measure the shortest route traps 2.8\% of workers when the warning comes 30 minutes late, which is 13 times the 0.2\% who would have no way out even with the exact source known.

\section{Routes checked against every source in the zone}\label{sec:fix}
The robust route protects nearly all outside starting points for a northwest alarm: 1,200 of 1,243 (96.5\%), up from 720. For a southwest alarm it protects 940 of 1,139, up from 659 (Table~\ref{tab:plans}). The fix usually means going to a different shaft. Of the 480 northwest starting points that only the robust plan protects, 462 end at a different shaft from their shortest route; in the southwest the figure is 207 of 281. Because it must work for every source in the zone, the robust route reaches the shaft cluster by an approach that none of those floods closes in time.

\begin{table}[htbp]
\centering\small
\caption{Starting points outside each quadrant zone that are protected at a 30-minute delay target. The robust columns allow routes up to 20\% or 50\% longer than the shortest route, or of any length. The last column counts starting points that the exact source would protect, the ceiling for any fixed plan.}\label{tab:plans}
\begin{tabular}{@{}lrrrrrrrr@{}}\toprule
& & & & & \multicolumn{3}{c}{Robust route} & \\
\cmidrule(lr){6-8}
Zone & Sources & Starting points & Shortest & Avoid zone & $\le$+20\% & $\le$+50\% & Any & Exact source\\\midrule
NW & 344 & 1,243 & 720 & 1,166 & 812 & 961 & 1,200 & 1,223\\
SW & 448 & 1,139 & 659 & 838 & 754 & 760 & 940 & 1,071\\
SE & 346 & 1,241 & 1,196 & 1,233 & 1,224 & 1,231 & 1,240 & 1,240\\
NE & 449 & 1,138 & 1,075 & 1,102 & 1,079 & 1,092 & 1,110 & 1,130\\
\bottomrule

\end{tabular}
\end{table}

The simple avoid-the-zone rule does surprisingly well on this count. It protects 1,166 northwest starting points and 838 southwest ones, so it recovers 446 of the 480 northwest starting points that the robust route adds and 179 of the 281 southwest gains (Figure~\ref{fig:plans}). Where it fails, it fails because it is blind to the flood: it treats every tunnel of the zone as equally dangerous and every tunnel outside it as safe.

The rule pays for its simplicity in walking. For the northwest starting points it newly protects, its median route is 364\% longer than the shortest route; the robust route protects the same workers with a median of 54\% extra. The robust route cannot be longer, since it is the shortest route with the guarantee, so the difference measures what the rule wastes. In the southwest, where the robust route itself is long (a median of 267\% extra for newly protected starting points), the two medians are the same. The detours that robust routes need are not always small either. In the northwest, a 20\% allowance protects 812 starting points and a 50\% allowance 961, against 1,200 with no limit (Table~\ref{tab:plans}). A practical plan would cap walking distance or time, and the same calculation shows how much protection remains under the cap.

\begin{figure}[tbp]
\centering\includegraphics[width=.86\linewidth]{figures/plans.pdf}
\caption{Share of starting points outside each quadrant zone protected by each plan at a 30-minute delay target. The gray bar spans the shortest route and the ceiling set by the exact source.}\label{fig:plans}
\end{figure}

Figure~\ref{fig:example} shows a typical case. A worker starts in the far northeast, well outside the northwest zone. The shortest route, 6.88\,km along the main roadway, ends at shaft P0986 at the northern end of the shaft cluster, and a flood from the northwest zone can cut it 130\,m before the shaft even if the worker leaves at once. The robust route runs down a parallel roadway to the southernmost shaft, P1081. It is 8.36\,km long, 22\% more, and it works for every northwest source as long as the worker sets off within about four hours of the inrush.

\begin{figure}[tbp]
\centering\includegraphics[width=.8\linewidth]{figures/example.pdf}
\caption{A worker starting at junction P0014 under a northwest alarm. The shortest route (red) follows the main roadway to shaft P0986 and can be cut near the shaft (black cross) by a flood from the northwest zone (purple dots), even with no delay. The robust route (blue) reaches shaft P1081 along a parallel roadway and is safe for every northwest source.}\label{fig:example}
\end{figure}

Workers inside the alarm zone are much harder to protect, because the source may be next to them. Robust routes raise the number protected from 14 to 97 of the 344 northwest starting points, from 31 to 139 of 448 in the southwest, from 25 to 138 of 346 in the southeast, and from 28 to 179 of 449 in the northeast (Appendix~\ref{app:inside}). The avoid-the-zone rule helps little here and in the southwest protects fewer workers than the shortest route, since a worker inside the zone cannot avoid it.

Over the whole mine, robust routes raise the guaranteed share of (worker, flood) pairs from 58.7\% to 78.8\%, and the avoid-the-zone rule to 69.2\%. If the robust plan falls back on the avoid-the-zone rule where no robust route exists, it traps 1.3\% of workers at a 30-minute delay, compared with 2.8\% for the shortest route and 1.8\% for the rule alone.

\section{Information or time?}\label{sec:info}
\subsection{The gap is a matter of information}
Under quadrant alarms, 21.2\% of (worker, flood) pairs remain unprotected by robust routes at a 30-minute target. Only 0.2\% have no way out, in the sense that even the exact source would not save the worker. The other 21.0\% are information-limited: a route exists for that flood, but no single route works for every source in the zone. Almost the whole gap is therefore a matter of information.

\subsection{Finer zones, and where their boundaries fall}
To see how much information is needed, we replaced the quadrants by zones of equal tunnel length, from a single zone covering the whole mine up to 128 zones (Table~\ref{tab:resolution} and Figure~\ref{fig:information}a). With no location information at all, robust routes can be guaranteed for only 10\% of pairs. Protection then rises with the number of zones, to 51--61\% with four zones, 82\% with sixteen, and 97\% with sixty-four, close to the 99.8\% ceiling. The three cutting angles agree within a few points, except with four, eight and thirty-two zones, where a small shift of the boundaries changes the result by up to nine points.

\begin{table}[htbp]
\centering\small
\caption{Share of (worker, flood) pairs (\%) protected at a 30-minute delay target, for alarm zones of equal tunnel length. Where the three cutting angles differ, the mean is followed by the range. The last column is the information-limited share: pairs that the exact source would protect but the zone does not.}\label{tab:resolution}
\begin{tabular}{@{}lllll@{}}\toprule
Zones & Shortest route & Avoid zone & Robust route & Information-limited\\\midrule
1 & 2.3 & 2.3 & 10.1 & 89.7\\
2 & 23.9 (22.4--25.1) & 23.8 (23.2--25.0) & 33.6 (31.1--36.0) & 66.1 (63.8--68.7)\\
4 & 42.2 (36.7--49.1) & 44.9 (39.1--53.1) & 55.4 (51.3--60.5) & 44.4 (39.3--48.5)\\
8 & 58.9 (55.5--64.1) & 61.7 (59.4--66.2) & 72.9 (69.1--76.7) & 26.9 (23.1--30.7)\\
16 & 68.9 (68.0--70.6) & 70.7 (69.1--73.3) & 81.8 (81.5--82.0) & 18.0 (17.8--18.3)\\
32 & 76.8 (74.3--80.9) & 82.1 (77.1--89.1) & 92.2 (87.1--95.0) & 7.5 (4.8--12.7)\\
64 & 84.6 (83.0--87.1) & 90.3 (87.0--92.2) & 97.2 (97.1--97.4) & 2.6 (2.4--2.7)\\
128 & 87.8 (86.2--89.6) & 92.8 (89.7--95.4) & 98.3 (98.1--98.5) & 1.5 (1.3--1.7)\\
\midrule Exact source & 97.2 & -- & 99.8 & 0.0\\
\bottomrule

\end{tabular}
\end{table}

The four quadrants stand out. They protect 78.8\% of pairs, about as much as sixteen equal zones and far more than four. The reason is again the shafts. In every equal-length partition with up to sixteen zones, a single zone contains all 167 junctions within 300\,m of a shaft. When that zone raises the alarm, any source in it can flood the approach to every shaft, and robust routes exist for only 10--33\% of starting points. The large step between sixteen and sixty-four zones is the point at which the boundaries first split the shaft area. The quadrant boundaries happen to cross right at the shafts, splitting this area almost in half, so that even their worst alarm leaves 82\% of starting points protected.

To test this directly, we drew zones through the shaft cluster on purpose. One straight line across the main roadway between the shafts gives two zones with two shafts on each side; a second line along the roadway gives four zones. Two zones divided at the shafts protect 51.7\% of pairs, against 31--36\% for two zones of equal tunnel length, and four protect 69.0\%, against 51--61\% (Figure~\ref{fig:information}a). Drawing the boundaries through the shafts is worth about as much as doubling the number of zones. It is not the whole story, since the quadrants, which also split the shaft area, do better still; where the other boundaries fall matters too.

\begin{figure}[tbp]
\centering\includegraphics[width=\linewidth]{figures/information.pdf}
\caption{(a) Share of (worker, flood) pairs protected at a 30-minute target, against the number of equal-length alarm zones; bands span the three cutting angles. The star marks the four quadrants and the diamonds the zones divided at the shafts. The gray dot at the right is the shortest route when the exact source is known. The dashed line is the ceiling set by the exact source. (b) The same share for quadrant zones, against the warning delay. Knowing the source makes delay almost harmless; with zone alarms, the plan loses protection with every minute of delay.}\label{fig:information}
\end{figure}

\subsection{Delay makes information more valuable}
Delay by itself costs very little in this model. With the exact source known, the share of pairs with a way out falls only from 100\% to 99.8\% as the delay grows from zero to 30 minutes, and to 99.5\% at two hours (Figure~\ref{fig:information}b). Under quadrant alarms, the picture is different. Robust routes protect 96.9\% of pairs when workers leave at once, 84.1\% at 10 minutes, 78.8\% at 30 minutes and 69.4\% at two hours. The information-limited share grows accordingly, from 3\% to 30\%.

Delay and information therefore interact. A worker who leaves immediately can almost always find one route that stays ahead of every possible flood in the zone. The longer the warning takes, the more the floods from different sources have spread, and the less likely it is that one route avoids all of them. With finer zones the effect shrinks: with 128 zones, protection falls only from 98.9\% at 10 minutes to 97.9\% at 60 minutes. A faster alarm and a more precise alarm are partly substitutes, but the precise one is worth more.

\section{Sensitivity to inflow and to the set of possible sources}\label{sec:robust}
\paragraph{A stronger inflow.}
Doubling the inflow to $1.0\,\mathrm{m^3/s}$ halves every flood time. With quadrant alarms and a 30-minute target, the guaranteed share falls to 56.9\% for shortest routes, 66.5\% for the avoid-the-zone rule and 71.6\% for robust routes, while the ceiling stays at 99.7\%. The pattern across zone resolutions is unchanged, including the jump at the point where the shaft area is split (Appendix~\ref{app:extra}). Because flood times scale with the inflow (Appendix~\ref{app:flood}) and walking is slower in wet tunnels, in this model a route that works at a higher inflow also works at any lower one, so a plan designed for an upper bound on the inflow covers all smaller inflows.

\paragraph{Fewer possible sources.}
So far every junction in the reported zone could be the source. In practice some places are much more likely than others. As a contrasting case, we kept only junctions within a given distance of the water areas marked on the CAD plan, which we treat as a rough stand-in for likely sources (Table~\ref{tab:breadth}). With sources within 50\,m of these areas, 211 junctions in all, shortest routes can already be guaranteed for 93.8\% of pairs, robust routes raise this to 97.8\%, and the information-limited share is 2.1\%. As the radius grows to 1\,km, the information-limited share rises to 19.0\%, close to the 21.0\% for whole zones. The fewer places the water can come from, the less the alarm needs to say, and the less a robust plan adds. Any estimate of what rerouting is worth therefore depends on which sources the plan is meant to withstand.

\begin{table}[htbp]
\centering\small
\caption{Share of (worker, flood) pairs (\%) protected under quadrant alarms at a 30-minute target, when the possible sources are restricted to junctions within a given distance of the water areas marked on the CAD plan.}\label{tab:breadth}
\begin{tabular}{@{}lrrrrrr@{}}\toprule
Possible sources & Junctions & Shortest & Avoid zone & Robust & Information-limited & Exact source\\\midrule
Within 50\,m & 211 & 93.8 & 93.7 & 97.8 & 2.1 & 99.9\\
Within 100\,m & 258 & 92.4 & 92.6 & 96.8 & 3.1 & 99.9\\
Within 200\,m & 436 & 82.7 & 87.1 & 93.3 & 6.5 & 99.8\\
Within 500\,m & 871 & 70.6 & 78.0 & 87.9 & 11.9 & 99.8\\
Within 1000\,m & 1,346 & 60.1 & 70.7 & 80.8 & 19.0 & 99.8\\
Whole zone & 1,587 & 58.7 & 69.2 & 78.8 & 21.0 & 99.8\\
\bottomrule

\end{tabular}
\end{table}

\section{Limitations and implications}\label{sec:limits}
The results come from a single mine, so they show what can happen rather than how often it happens in mines generally. In particular, the strong effect of the shafts depends on this mine's layout, with all four shafts close together on one roadway; a mine with shafts far apart would behave differently. Our zones are geometric stand-ins for sensor zones. We treated every junction as an equally likely source and did not consider points partway along a tunnel. The counts in our tables are numbers of junctions, which are not the same as numbers of workers.

The flood model is the main physical limitation. It conserves water volume, but it splits flow equally at junctions, spreads water along a tree of tunnels, and tracks only the mean depth of each tunnel. It ignores momentum, changes in depth along a tunnel, losses at junctions, pumping, debris, and the force of moving water on a person. Its timing has not been calibrated against any measurement. The extracted geometry and estimated elevations would also need review by mining engineers.

The model of the workers is equally simple. Routes end at the bottom of a shaft, and the model ignores queues, hoist capacity, fatigue, injury, and the chance of a wrong turn. Some robust routes are long, and a long route may be safe in the model and still be a poor emergency instruction.

Within these limits, the study suggests three practical points. First, check the shortest routes against every source in each alarm zone, not just against the most likely one; a route can fail far from where the water starts. Second, look at the area around the shafts: if all shaft approaches lie in one alarm zone, that zone is the weak point of the plan, and splitting it is worth more than adding zones elsewhere. Third, a simple rule of avoiding the alarm zone captures much of the benefit of rerouting, but checking routes against a flood model can find much shorter ones and can show which workers no fixed route can protect.

\section{Conclusion}
A fixed escape plan has to work without knowing where the water comes from. In the mine we studied, the shortest routes fail for many workers far from any flood, because the water reaches the shaft area first. Routes checked against every possible source in the alarm zone repair most of these failures, and a simple rule of avoiding the zone repairs many of them at a higher walking cost. What remains is almost entirely a matter of information: alarm zones that narrow down the source, and in particular zones that separate the approaches to different shafts, would let a fixed plan protect nearly every worker, while an earlier warning helps mainly by making that information less critical. The natural next steps are to test these findings with a calibrated flood model and with mines whose shafts are spread out, and to design alarm zones and escape routes together.

\begin{thebibliography}{99}\small
\bibitem{bayram2016} V. Bayram. Optimization models for large scale network evacuation planning and management: A literature review. \emph{Surveys in Operations Research and Management Science}, 21(2):63--84, 2016. \href{https://doi.org/10.1016/j.sorms.2016.11.001}{doi:10.1016/j.sorms.2016.11.001}.
\bibitem{bental2009} A. Ben-Tal, L. El Ghaoui, and A. Nemirovski. \emph{Robust Optimization}. Princeton University Press, 2009. \href{https://doi.org/10.1515/9781400831050}{doi:10.1515/9781400831050}.
\bibitem{bertsimas2003} D. Bertsimas and M. Sim. Robust discrete optimization and network flows. \emph{Mathematical Programming}, 98(1--3):49--71, 2003. \href{https://doi.org/10.1007/s10107-003-0396-4}{doi:10.1007/s10107-003-0396-4}.
\bibitem{cumcm} China Society for Industrial and Applied Mathematics. 2025 Higher Education Press Cup China Undergraduate Mathematical Contest in Modeling (CUMCM), Problem D: Spreading model of mine water inrush and escape plan. Problem statement and network data, 2025.
\bibitem{dreyfus1969} S. E. Dreyfus. An appraisal of some shortest-path algorithms. \emph{Operations Research}, 17(3):395--412, 1969. \href{https://doi.org/10.1287/opre.17.3.395}{doi:10.1287/opre.17.3.395}.
\bibitem{hamacher2006} H. W. Hamacher, S. Ruzika, and S. A. Tjandra. Algorithms for time-dependent bicriteria shortest path problems. \emph{Discrete Optimization}, 3(3):238--254, 2006. \href{https://doi.org/10.1016/j.disopt.2006.05.006}{doi:10.1016/j.disopt.2006.05.006}.
\bibitem{hanasusanto} G. A. Hanasusanto, D. Kuhn, and W. Wiesemann. $K$-adaptability in two-stage robust binary programming. \emph{Operations Research}, 63(4):877--891, 2015. \href{https://doi.org/10.1287/opre.2015.1392}{doi:10.1287/opre.2015.1392}.
\bibitem{kaufman} D. E. Kaufman and R. L. Smith. Fastest paths in time-dependent networks for intelligent vehicle-highway systems application. \emph{IVHS Journal}, 1(1):1--11, 1993. \href{https://doi.org/10.1080/10248079308903779}{doi:10.1080/10248079308903779}.
\bibitem{martins1984} E. Q. V. Martins. On a multicriteria shortest path problem. \emph{European Journal of Operational Research}, 16(2):236--245, 1984. \href{https://doi.org/10.1016/0377-2217(84)90077-8}{doi:10.1016/0377-2217(84)90077-8}.
\bibitem{millerhooks2000} E. D. Miller-Hooks and H. S. Mahmassani. Least expected time paths in stochastic, time-varying transportation networks. \emph{Transportation Science}, 34(2):198--215, 2000. \href{https://doi.org/10.1287/trsc.34.2.198.12304}{doi:10.1287/trsc.34.2.198.12304}.
\bibitem{mutapcic2009} A. Mutapcic and S. Boyd. Cutting-set methods for robust convex optimization with pessimizing oracles. \emph{Optimization Methods and Software}, 24(3):381--406, 2009. \href{https://doi.org/10.1080/10556780802712889}{doi:10.1080/10556780802712889}.
\bibitem{nrc2013} National Research Council. \emph{Improving Self-Escape from Underground Coal Mines}. The National Academies Press, Washington, DC, 2013. \href{https://doi.org/10.17226/18300}{doi:10.17226/18300}.
\bibitem{orda} A. Orda and R. Rom. Shortest-path and minimum-delay algorithms in networks with time-dependent edge-length. \emph{Journal of the ACM}, 37(3):607--625, 1990. \href{https://doi.org/10.1145/79147.214078}{doi:10.1145/79147.214078}.
\bibitem{sun2016} W. Sun, W. Zhou, and J. Jiao. Hydrogeological classification and water inrush accidents in China's coal mines. \emph{Mine Water and the Environment}, 35(2):214--220, 2016. \href{https://doi.org/10.1007/s10230-015-0363-3}{doi:10.1007/s10230-015-0363-3}.
\bibitem{sun2024} W. Sun, W. Li, L. Ren, and K. Li. Spatial and temporal characterization of mine water inrush accidents in China, 2014--2022. \emph{Water}, 16(5):656, 2024. \href{https://doi.org/10.3390/w16050656}{doi:10.3390/w16050656}.
\bibitem{wu2020} Q. Wu, Y. Du, H. Xu, Y. Zhao, X. Zhang, and Y. Yao. Finding the earliest arrival path through a time-varying network for evacuation planning of mine water inrush. \emph{Safety Science}, 130:104836, 2020. \href{https://doi.org/10.1016/j.ssci.2020.104836}{doi:10.1016/j.ssci.2020.104836}.
\bibitem{wu2021} Q. Wu, Z. Du, Y. Zhao, H. Xu, and X. Zhang. Optimal location of water level sensors for monitoring mine water inrush based on the set covering model. \emph{Scientific Reports}, 11:2621, 2021. \href{https://doi.org/10.1038/s41598-021-82121-7}{doi:10.1038/s41598-021-82121-7}.
\bibitem{yu1998} G. Yu and J. Yang. On the robust shortest path problem. \emph{Computers \& Operations Research}, 25(6):457--468, 1998. \href{https://doi.org/10.1016/S0305-0548(97)00085-3}{doi:10.1016/S0305-0548(97)00085-3}.
\bibitem{zhang2022} X. Zhang, Q. Wu, Y. Zhao, S. Liu, and H. Xu. Spatiotemporal modeling of water inrush spreading in mine roadway networks. \emph{Water Science and Technology}, 85(3):872--886, 2022. \href{https://doi.org/10.2166/wst.2022.009}{doi:10.2166/wst.2022.009}.
\end{thebibliography}
\appendix
\section{Travel times, guarantees, and exact search}\label{app:math}
Represent each tunnel by two directed arcs. For arc $a$ in the flood from source $s$, let $\ell_a>0$ be its length, $w_a^s$ the time it gets wet, $b_a^s$ the time it closes, and $v_a^s$ the walking speed once it is wet. The dry speed satisfies $v_d\ge v_a^s>0$. Ignoring closure, a worker who enters the arc at time $t$ arrives at its far end at
\[
 \widetilde A_a^s(t)=
 \begin{cases}
 t+\ell_a/v_d,&\ell_a\le v_d(w_a^s-t)_+,\\
 \max(t,w_a^s)+[\ell_a-v_d(w_a^s-t)_+]/v_a^s,&\text{otherwise},
 \end{cases}
\]
where $(x)_+=\max(x,0)$. Set $A_a^s(t)=\widetilde A_a^s(t)$ if this is at most $b_a^s$, and $A_a^s(t)=+\infty$ otherwise. Arriving exactly at the closure time counts as success; a stricter safety convention would add an explicit margin.

For a path $P=(a_1,\ldots,a_m)$, composing the arc functions gives $A_P^s(t)$. These functions are nondecreasing (counting $+\infty$ as infeasible), so traversing a path earlier cannot turn a feasible journey into an infeasible one. The latest safe departure along $P$ is
\[
 D_s(P)=\sup\{d\ge0:A_P^s(d)<\infty\},\qquad
 D_Z(P)=\min_{s\in\Sc_Z}D_s(P),
\]
with $D_s(P)=-\infty$ if no nonnegative departure works.

\begin{proposition}[Checking one departure time is enough]
For the permanent-closure travel times above, a path $P$ works in every flood $s\in\Sc_Z$ for every departure $d\in[0,\tau]$ if and only if $A_P^s(\tau)<\infty$ for every $s\in\Sc_Z$.
\end{proposition}
\begin{proof}
Necessity follows by taking $d=\tau$. For sufficiency, monotonicity gives $A_P^s(d)\le A_P^s(\tau)<\infty$ for every earlier $d$ and every flood.
\end{proof}

\paragraph{Backward deadlines.}
For a tunnel with parameters $(\ell,w,b,v_d,v_w)$ and a deadline $q$ at its far end, set $e=\min(q,b)$. The latest entry time that reaches the far end by $e$ is
\[
 B(q)=\begin{cases}
 e-\ell/v_w,&\ell\le v_w(e-w)_+,\\
 \min(e,w)-[\ell-v_w(e-w)_+]/v_d,&\text{otherwise}.
 \end{cases}
\]
Applying this map from the last arc back to the first evaluates a fixed path. A reverse label-setting search maximizes it over all routes for one known source; call the result $U_s(v)$. The exact-source plan protects the pair $(v,s)$ exactly when $U_s(v)\ge\tau$; a pair with $U_s(v)<\tau$ has no way out.

\begin{proposition}[Adding sources one at a time]
Fix a starting point $v$, a target $\tau$, and a finite source set $\Sc_Z$. Suppose each search either returns an exact shortest route that works for its active subset $T\subseteq\Sc_Z$, or certifies that none exists. If a returned route works for all of $\Sc_Z$, its length is $L_Z^*(v,\tau)$. If the search for some subset finds no route, no route works for $\Sc_Z$. Adding a violated source after every failed check reaches one of these conclusions after finitely many additions.
\end{proposition}
\begin{proof}
Every route that works for $\Sc_Z$ also works for $T$, so the subset optimum is a lower bound on the full optimum. If the subset route passes all sources, it is also an upper bound, and the two agree. Infeasibility for $T$ implies infeasibility for any larger set. A violated source cannot already be in $T$, so each failed check strictly enlarges $T$, which stays inside the finite set $\Sc_Z$.
\end{proof}
Floating-point feasibility checks use tolerances of $10^{-7}$ seconds or less, which are not a physical safety margin.

\paragraph{Two consequences used in the paper.}
First, splitting a zone can only help: a route that works for every source in $Z$ works for every source in any part of $Z$, so the robust plan for a finer partition protects every pair that the coarser one protects. Second, a menu of routes indexed by the known delay cannot extend the guarantee for a fixed zone: to cover all delays up to $\tau$, the menu must include, for delay $\tau$, a route that works for every source, and by the first proposition that route alone also works for every earlier delay.

\section{Flood propagation rules}\label{app:flood}
Each tunnel stores a volume between zero and its capacity $WH\ell$, where $W=4.99$\,m and $H=3$\,m. When water first reaches a level or downhill tunnel, the tunnel receives a layer of volume $Wh_0\ell$ with $h_0=0.1$\,m. If its far junction is still dry, the tunnel then passes water on and joins a propagation tree; if the far junction is already wet, it stores further water. An uphill tunnel fills from below before passing water to a dry junction above it. At a wet junction, incoming flow is divided equally among the level or downhill tunnels still available (with a 0.005\,m elevation tolerance), and uphill tunnels are used only when none remain. When a dead end fills up, the tunnel feeding it stops passing water on and stores it instead. Inflow trapped behind full tunnels is moved through them to the nearest junction that can still store water, and overflow after saturation is counted in the volume balance. Simultaneous events are ordered by junction and tunnel identifiers. A tunnel closes when its mean depth reaches 0.3\,m. Walking switches to wet speed when water first enters a tunnel; the model does not track a moving front inside a tunnel.

Multiplying a constant inflow by a factor divides every flood event time by the same factor, so the doubled-inflow floods are the nominal floods with all times halved. A direct rerun for one source reproduced the scaled times exactly.

\section{Checks on the implementation}\label{app:checks}
We checked the results in four ways.
\begin{enumerate}
\item \emph{Brute force on small networks.} The search code was tested on 100 small random networks with three floods each, against enumeration of every simple path; all 400 comparisons of feasibility and minimum length agreed.
\item \emph{Reproduction.} For the quadrant zones at 10, 20 and 30 minutes, the current code reproduces an earlier run of the same 19,044 queries, made with separate driver code, exactly, including every route length.
\item \emph{Logical consistency.} Several relations must hold if the search is exact. The robust route is the shortest protected route, so it must exist whenever the shortest route or the avoid-the-zone route is protected, and it can be no longer than either. It can exist only if every source in the zone leaves a way out. Its latest safe departure must reach the target. And because a route that works for a zone also works for each half of it, every starting point protected for a zone must be protected for both halves in the next finer partition. All 2,558,244 records satisfy the first four relations, and all 747 parent--child zone pairs satisfy the last.
\item \emph{Independent forward walks.} For each quadrant zone we sampled eight robust and eight zone-avoiding routes and walked each one forward through every flood in its zone, with code separate from the search. Leaving just before the computed latest departure reached a shaft in all 23,012 route--source checks, and leaving one second after it failed for at least one source in every case.
\end{enumerate}
The flood simulations conserve water volume to a relative error of at most $1.02\times10^{-13}$. This verifies numerical mass conservation.

\section{Workers inside the alarm zone}\label{app:inside}
\begin{table}[htbp]
\centering\small
\caption{Starting points inside each quadrant zone at a 30-minute target. ``No way out'' counts starting points for which at least one source in the zone leaves no route even if the source is known; ``No common route'' counts those for which every source leaves a route, but no single route works for all of them.}
\begin{tabular}{@{}lrrrrrr@{}}\toprule
Zone & Starting points & Shortest & Avoid zone & Robust & No way out & No common route\\\midrule
NW & 344 & 14 & 33 & 97 & 172 & 75\\
SW & 448 & 31 & 21 & 139 & 202 & 107\\
SE & 346 & 25 & 55 & 138 & 105 & 103\\
NE & 449 & 28 & 13 & 179 & 215 & 55\\
\bottomrule

\end{tabular}
\end{table}

\section{Delay, inflow and zone resolution}\label{app:extra}
\begin{table}[htbp]
\centering\small
\caption{Share of (worker, flood) pairs (\%) guaranteed by robust routes for several alarm partitions, at three delay targets with the nominal inflow of $0.5\,\mathrm{m^3/s}$, and at 30 minutes with the inflow doubled. The equal-length zones use the first cutting angle.}\label{tab:extra}
\begin{tabular}{@{}lrrrr@{}}\toprule
& \multicolumn{3}{c}{$0.5\,\mathrm{m^3/s}$} & $1.0\,\mathrm{m^3/s}$\\\cmidrule(lr){2-4}\cmidrule(lr){5-5}
Alarm zones & 10 min & 30 min & 60 min & 30 min\\\midrule
1 equal zone (no location) & 16.4 & 10.1 & 3.5 & 1.4\\
2 equal zones & 43.2 & 36.0 & 28.3 & 25.6\\
4 equal zones & 59.9 & 54.5 & 48.7 & 47.4\\
8 equal zones & 76.7 & 72.9 & 67.2 & 63.8\\
16 equal zones & 84.8 & 82.0 & 77.0 & 74.0\\
32 equal zones & 91.2 & 87.1 & 81.2 & 79.6\\
64 equal zones & 98.2 & 97.4 & 96.6 & 95.9\\
128 equal zones & 98.9 & 98.4 & 97.9 & 97.6\\
\midrule
4 quadrants & 84.1 & 78.8 & 73.6 & 71.6\\
2 zones cut at the shafts & 54.8 & 51.7 & 49.1 & --\\
4 zones cut at the shafts & 73.2 & 69.0 & 63.5 & --\\
Exact source & 99.9 & 99.8 & 99.7 & 99.7\\
\bottomrule

\end{tabular}
\end{table}
\end{document}
